\chapter{Background and Technologies}
\label{chap:background}

\section{GPU Sharing}
\label{sec:gpu-sharing}

\begin{itemize}
    \item Describe the available GPU-sharing mechanisms and compare their isolation, flexibility, and resource-efficiency characteristics.
    \item Explain why NVIDIA MIG was selected for the experimental work presented in this thesis.
\end{itemize}

\section{NVIDIA Multi-Instance GPU}
\label{sec:mig}

\begin{itemize}
    \item Provide an in-depth description of the MIG architecture, supported instance profiles, resource partitioning, isolation properties, configuration workflow, and relevant limitations.
\end{itemize}

\section{GPUs in Kubernetes}
\label{sec:gpus-kubernetes}

\begin{itemize}
    \item Explain how Kubernetes traditionally discovers, advertises, schedules, and exposes GPU resources through device plugins.
    \item Discuss the implications of extended resources and device-plugin-based allocation for GPU workloads.
\end{itemize}

\section{Dynamic Resource Allocation}
\label{sec:dra}

\begin{itemize}
    \item Provide an in-depth description of the DRA architecture and its principal API objects.
    \item Explain the resource claim lifecycle, device selection, scheduling integration, and allocation semantics relevant to NVIDIA GPUs and MIG devices.
    \item Clarify the differences between DRA and conventional device-plugin-based allocation, including current operational constraints.
\end{itemize}

\section{Linux GPU Software Stack}
\label{sec:linux-gpu-stack}

\begin{itemize}
    \item Describe the Linux software stack required to expose NVIDIA GPUs to containers, including the driver, container runtime integration, device discovery, and monitoring components.
\end{itemize}

\section{NVIDIA GPU Operator}
\label{sec:gpu-operator}

\begin{itemize}
    \item Explain how the NVIDIA GPU Operator deploys and manages the software components required for GPU-enabled Kubernetes clusters.
    \item Describe its role in configuring MIG and enabling NVIDIA's DRA driver in the experimental environment.
\end{itemize}

\section{Architecture Overview}
\label{sec:architecture-overview}

\begin{itemize}
    \item Organize the system into three conceptual layers: the hardware layer, the resource-allocation layer, and the application layer.
    \item Describe the responsibilities of each layer and the interactions among MIG configuration, Kubernetes resource allocation, and AI inference workloads.
\end{itemize}
